\documentclass[11pt, a4paper]{article}
\usepackage{graphicx} % Required for inserting images
\usepackage{natbib}
\usepackage{amsmath}
\usepackage{hyperref}
\usepackage[left=1.5cm, right=2cm]{geometry}
\title{Detectability of 4FGL catalog \textit{Fermi}-LAT sources with Imaging Atmospheric Cherenkov Telescopes}
\author{Sandra Dembińska$^1$, Wojciech Jabłoński$^1$, Kajetan Janusz $^1$, Mateusz Jaworski$^1$,\\
Bartosz Palewicz$^1$, Jakub Szczech$^1$, Krzysztof Wilk$^1$, Julian Sitarek$^1$\\
\\
$^1$Faculty of Physics and Applied Informatics, University of Lodz
}
\date{April 2026}

\begin{document}

\maketitle

\section{Introduction}
The purpose of the project is to prepare an application that aggregates information from various sources in order to evaluate detectability of the sources listed in the 4FGL-DR4 catalog \citep{2020ApJS..247...33A}.
The application includes multi-TeV information from HAWC and LHAASO experiments, as well as various multiwavelength data available in NASA NED database.
The sources are cross-references according to the known TeV sources in TeVCAT.
In order to evaluate detectability of the sources MAGIC Source Simulator (mss) is used.

\section{Project concept and task management, MVP -- Jakub Szczech}

Jakub Szczech was responsible for defining the overall project concept and translating it into a concrete development roadmap.
The project addresses the need for a unified, web-accessible interface that aggregates high-energy gamma-ray catalogs — \textit{Fermi}-LAT 4FGL, HAWC 3HWC, LHAASO, and TeVCat — alongside multiwavelength data from the NASA NED database, enabling astronomers to assess the detectability of known sources with Imaging Atmospheric Cherenkov Telescopes (IACTs).

The minimum viable product (MVP) was scoped to include four essential components:
\begin{enumerate}
    \item A catalog ingestion pipeline capable of downloading, parsing, and cross-matching sources from at least the \textit{Fermi}-LAT and HAWC catalogs.
    \item A PostGIS-backed relational data model storing unified source records together with catalog-specific metadata.
    \item A REST API exposing paginated source listings, spatial cone searches, and advanced multi-field filtering.
    \item A browser-based single-page application allowing users to browse, filter, and visualise the aggregated source set.
\end{enumerate}

Each MVP component was treated as an independent track — backend data model and ingestion, REST API, and frontend — so that parallel development could proceed without blocking dependencies.
Tasks were decomposed into discrete work items with defined acceptance criteria specifying the expected API contract, data model constraints, and UI behaviour.
This structured approach ensured that integration points between components (for example, the shape of the JSON envelope returned by paginated endpoints) were agreed upon before implementation, reducing friction during final integration.

Feature prioritisation followed the principle that every item in the MVP must be demonstrably correct before optional extensions — such as the 3D celestial sphere, analytics dashboard, or MAGIC Source Simulator integration — were begun.
This sequencing allowed the team to maintain a working, deployable system at every stage of development, rather than accumulating unintegrated work across tracks.

The project management strategy proved effective in keeping individual contributions focused and testable.
Future extensions, such as real-time ingestion of newly released catalog versions or integration with additional IACT simulation frameworks, can be incorporated by adding new loader classes and API endpoint actions without altering the unified source model.

\section{Repository management and integration - Kajetan Janusz}

Kajetan Janusz managed the repository and was responsible for the full containerisation and continuous integration infrastructure.
The goal of this work was to provide every contributor with a reproducible, single-command local development environment and to enforce code quality gates on every pull request through automated pipelines.

The application is orchestrated with Docker Compose, defining three services that form a strict health-check dependency chain:

\begin{enumerate}
    \item \textbf{database} — PostgreSQL~16 with PostGIS~3.4 (\texttt{postgis/postgis:16-3.4-alpine}), exposed on port~5432, with a named volume \texttt{postgres-data} for persistence.
    \item \textbf{backend} — Django REST API built from \texttt{backend/Dockerfile}, exposed on port~8000; starts only after \texttt{database} passes its \texttt{pg\_isready} health check.
    \item \textbf{frontend} — Vite development server built from \texttt{frontend/Dockerfile}, exposed on port~3000; starts only after \texttt{backend} passes its HTTP health check.
\end{enumerate}

The dependency ordering is expressed as:
\begin{equation}
    \texttt{database} \xrightarrow{\text{healthy}} \texttt{backend} \xrightarrow{\text{healthy}} \texttt{frontend}
\end{equation}

The backend health check polls the live catalog endpoint:
\begin{verbatim}
curl -fsS http://localhost:8000/api/sources/?page=1&page_size=1
\end{verbatim}
with a 180-second \texttt{start\_period} to accommodate the initial catalog download and ingestion, which can take several minutes on a cold start.

Separate Dockerfiles were written for each service.
The backend image is based on \texttt{python:3.12.13-slim-bookworm}; it installs PostGIS system libraries (\texttt{libgdal-dev}, \texttt{libgeos-dev}, \texttt{libproj-dev}), build toolchains, and Python dependencies via Poetry with \texttt{POETRY\_VIRTUALENVS\_CREATE=false} so packages are installed system-wide, avoiding virtualenv overhead in the container.
Django static files are collected at image build time via \texttt{manage.py collectstatic}.
The frontend image is based on \texttt{node:22.22.2-bookworm-slim}; it installs exact dependencies from \texttt{package-lock.json} via \texttt{npm ci} and runs the Vite development server bound to all interfaces on port~3000.

Two GitHub Actions workflows enforce quality on every pull request.
The \emph{backend} workflow triggers on changes to \texttt{backend/**} and runs two sequential jobs: (1) Ruff static analysis, and (2) pytest with PostGIS system libraries pre-installed.
The \emph{frontend} workflow triggers on changes to \texttt{frontend/**} and runs three independent parallel jobs: (1) ESLint, (2) TypeScript type-checking via \texttt{tsc --noEmit}, and (3) Jest unit tests.
Parallelism in the frontend pipeline reduces wall-clock time by executing all three checks simultaneously.

All environment variables are supplied via a single \texttt{.env} file at the repository root, referenced by both the backend and frontend services.
An \texttt{.env.example} file is committed to the repository documenting the required variables and their default values, enabling any contributor to start the full stack with \texttt{docker compose up -{}-build} after a single \texttt{cp .env.example .env}.

\section{Data aggregation from \textit{Fermi}-LAT, HAWC and LHAASO catalogs -- Bartosz Palewicz}

Bartosz Palewicz implemented the catalog ingestion pipeline that populates the unified source database from the \textit{Fermi}-LAT 4FGL, HAWC 3HWC, LHAASO, and TeVCat catalogs.
The pipeline is designed around a pluggable loader architecture: each catalog is handled by a dedicated class that encapsulates download, parsing, and normalisation logic, while a shared \texttt{CrossMatchService} handles spatial association with existing records.

Each catalog loader follows a common interface and the ingestion process proceeds in the following steps:

\begin{enumerate}
    \item \textbf{Download.} The loader fetches the catalog file (FITS or YAML) from its canonical upstream URL and caches it in the local \texttt{backend/data/} directory. Subsequent runs reuse the cached file.
    \item \textbf{Parse.} The file is read into memory. For FITS catalogs (\texttt{FermiLoader}, \texttt{LHASOLoader}) \texttt{astropy.table} reads binary table HDU~1. For the HAWC catalog (\texttt{HAWCLoader}) a YAML document is parsed with flexible column-name matching to handle variations between releases. TeVCat sources are retrieved via an astroquery TAP query against the HEASARC \texttt{tevcat} table.
    \item \textbf{Normalise.} Each row is converted to a common intermediate dictionary containing at minimum \texttt{ra}, \texttt{dec}, \texttt{name}, and a \texttt{metadata} payload preserving catalog-specific fields. Non-finite values are converted to \texttt{None}; bit-flag fields (e.g.\ \textit{Fermi} variability flags) are decoded into boolean attributes.
    \item \textbf{Cross-match.} The \texttt{CrossMatchService} queries the existing \texttt{Source} table for records within the match radius using a PostGIS spatial index.
    \item \textbf{Persist.} If a match is found the closest existing \texttt{Source} record is reused; a new \texttt{CatalogEntry} row is created (or updated) linking the incoming data to that source. If no match is found a new \texttt{Source} is created first.
\end{enumerate}

The \textit{Fermi}-LAT loader extracts the following fields from 4FGL-DR4 HDU~1: equatorial coordinates (\texttt{RAJ2000}, \texttt{DEJ2000}), source name and associated counterpart (\texttt{Source\_Name}, \texttt{ASSOC1}), source class (\texttt{CLASS1}), integrated energy flux above 1\,GeV (\texttt{Flux1000}), average significance (\texttt{Signif\_Avg}), test statistic (\texttt{Test\_Statistic}), spectral type and power-law index (\texttt{SpectrumType}, \texttt{PL\_Index}), redshift, and position uncertainty ellipse semi-axes.
The HAWC loader extracts flux at the decorrelation energy, spectral index, significance, and extension from the 3HWC YAML.
The LHAASO loader similarly extracts flux at 1\,TeV, spectral index, significance, and angular extension from the 1LHAASO FITS file.

The spatial cross-matching criterion is based on an angular proximity threshold $\theta_{\max}$.
The PostGIS \texttt{ST\_DWithin} function operates on the \texttt{position} field (SRID~4326, geographic coordinates), with the radius converted from degrees to kilometres using the approximation:
\begin{equation}
    d_{\max} \;[\text{km}] = \theta_{\max} \;[\text{deg}] \times 111
\end{equation}
The default match radius is $\theta_{\max} = 0.2^{\circ}$ ($\approx 12$\,arcmin).
When multiple existing sources fall within $\theta_{\max}$, the closest one by great-circle angular separation is selected.
The on-sky angular separation between a candidate position $(\alpha_1, \delta_1)$ and an existing source $(\alpha_2, \delta_2)$ is computed by PostGIS as:
\begin{equation}
    \theta = 2 \arcsin\!\left[\sqrt{\sin^2\!\frac{\delta_1 - \delta_2}{2} + \cos\delta_1\,\cos\delta_2\,\sin^2\!\frac{\alpha_1 - \alpha_2}{2}}\right]
\end{equation}

Catalog-specific metadata is preserved in a \texttt{JSONField} on the \texttt{CatalogEntry} model, allowing each catalog's native fields to be stored without schema migration when new catalogs are added.
The ingestion pipeline is exposed as a Django management command (\texttt{ingest\_catalogs}) that accepts a \texttt{-{}-catalogs} flag to ingest a subset of catalogs and a \texttt{-{}-match-radius} flag to override the default cross-matching threshold, facilitating sensitivity analysis of the association algorithm.

\section{UX/UI, documentation, user instructions -- Sandra Dembińska}

Sandra Dembińska was responsible for the project documentation, user experience design, and requirements specification.

Early in the project she prepared wireframes in Figma that defined the layout, component hierarchy, and navigation flow adopted by the frontend implementation.
The wireframes covered all four main views of the application, each corresponding to a top-level navigation entry: \emph{Strona główna}, \emph{Mapa wszechświata}, \emph{Analityka}, and \emph{Źródła}.

The \emph{Strona główna} (home page) wireframe establishes the global navigation bar shared across all views and presents an interactive 3D celestial sphere as the landing visual, with gamma-ray sources rendered as colour-coded points on the globe surface.
The \emph{Mapa wszechświata} wireframe places the 3D globe at the centre of the layout with two flanking side panels: a \emph{Legenda} panel on the left and a \emph{Filtry} panel on the right, both using checkbox lists to control source visibility per catalog.
The \emph{Analityka} wireframe defines a multi-section analytics dashboard: a catalog selector row (Fermi, HAWC, LHAASO, NED, TeVCat), a summary statistics table comparing catalogs across columns such as source count, average and median emission, average and peak significance, and high-detectability percentage, followed by KPI summary cards and a chart section containing an annual emission trend line chart, an average-vs-peak emission bar chart, and a statistical significance histogram.
The \emph{Źródła} wireframe specifies the source catalog browser: a filter panel exposing text search, source class selector, RA and Dec range inputs, minimum flux, and minimum significance fields, followed by a paginated data table with columns for source name, class, RA, Dec, Galactic longitude and latitude, flux, spectral index, and detection significance.

\begin{figure}[h]
    \centering
    \includegraphics[width=0.48\textwidth]{figures/wireframe_home.jpg}
    \hfill
    \includegraphics[width=0.48\textwidth]{figures/wireframe_universe_map.jpg}
    \caption{Wireframes for the \emph{Strona główna} (left) and \emph{Mapa wszechświata} (right) views. The home page presents the top navigation bar and an interactive 3D celestial sphere with colour-coded source points. The universe map view centres the same globe between a \emph{Legenda} panel (left) and a \emph{Filtry} panel (right), each containing per-catalog checkbox controls.}
    \label{fig:wireframes_map}
\end{figure}

\begin{figure}[h]
    \centering
    \includegraphics[width=0.48\textwidth]{figures/wireframe_analytics.jpg}
    \hfill
    \includegraphics[width=0.48\textwidth]{figures/wireframe_sources.jpg}
    \caption{Wireframes for the \emph{Analityka} (left) and \emph{Źródła} (right) views. The analytics dashboard shows a catalog selector, a per-catalog statistics table, KPI summary cards, and three chart sections (emission trend, average-vs-peak bar chart, significance histogram). The sources view presents a filter panel with coordinate range, flux, and significance inputs above a paginated data table listing source name, class, coordinates, flux, spectral index, and detection significance.}
    \label{fig:wireframes_data}
\end{figure}

She produced the full requirements specification for the application in collaboration with the entire team.
Requirements were derived iteratively through team discussions, with each contributor reviewing and refining the specification to ensure it accurately reflected both the scientific goals and the technical constraints of the system.
The final specification contains 66 functional requirements (FR-01–FR-66) and 16 non-functional requirements (NFR-01–NFR-16).
The functional requirements are grouped into seven areas: catalog ingestion (FR-01–FR-10), data model and storage (FR-11–FR-18), REST API (FR-19–FR-35), scoring and analytics calculations (FR-36–FR-39), the catalog browser frontend (FR-40–FR-45), the analytics dashboard frontend (FR-46–FR-57), the 3D Universe Map frontend (FR-58–FR-64), and internationalisation (FR-65–FR-66).
The non-functional requirements cover performance, startup and reliability, data persistence, spatial indexing, time and locale handling, accessibility, deployability, and pre-production security constraints.

Sandra Dembińska also produced the complete technical documentation suite.
The \emph{backend documentation} describes the full technology stack (Python~3.12, Django~6, Django REST Framework, PostGIS), the four-app Django project layout (\texttt{sources}, \texttt{catalogs}, \texttt{api}, \texttt{fermi}), the two core data models (\texttt{Source} and \texttt{CatalogEntry}), all seven REST API endpoints with their query parameters, and the catalog ingestion management commands for each of the four supported catalogs.
The \emph{frontend documentation} covers the technology stack (React~19, TypeScript~5.9, Vite~8, TanStack React Query, Three.js, D3, i18next), the full \texttt{src/} directory layout, the five-route React Router configuration, the three custom query hooks, the API module structure, Vite proxy configuration, environment variables, internationalisation setup, and test tooling.
The \emph{infrastructure documentation} details the three-service Docker Compose topology with its health-check dependency chain, volume and network configuration, both Dockerfiles (backend based on \texttt{python:3.12.13-slim-bookworm}, frontend on \texttt{node:22.22.2-bookworm-slim}), environment variable management via a single \texttt{.env} file, and common operational procedures.

Finally, she authored the \emph{user manual}, structured into eight chapters: Getting Started, Navigation, Home Page, Sources~--~Catalog Browser (with subsections on filters, the results table, and the source detail view), Analytics Dashboard (sky map, analysis scope, headline metrics, charts, and top sources table), Universe Map~--~3D View (controls, filters and legend, and the source detail panel), Data Reference (catalogs, source classes, and field definitions), and Language.
The manual is written to serve both the domain expert seeking specific catalog data and the general user unfamiliar with gamma-ray astronomy conventions.

\section{Application frontend  -- Wojciech Jabłoński}

Wojciech Jabłoński established the frontend architecture and built the shared infrastructure that all other frontend work depends on.
The application is a single-page application (SPA) built with React~19, TypeScript~5.9, and Vite~8, which provides sub-second hot-module replacement during development and optimised production bundles via Rollup.

The reusable component library was built as part of this contribution and provides a consistent foundation across all pages.
Key components include \texttt{DataTable} (a generic sortable table with configurable column definitions and loading/empty states), \texttt{DataTableFilters} (a popover-based filter panel supporting text, numeric, and multi-select filter types), \texttt{DataTablePagination} (page size and navigation controls), and \texttt{NavigationHeader} (responsive top bar with mobile hamburger menu, active-link highlighting, and language switcher).
Primitive UI elements — buttons, cards, badges, dialogs, sliders, and command menus — are sourced from the shadcn/ui library built on top of Radix UI, which provides accessible, unstyled components styled with Tailwind CSS~v4.
The layout is fully responsive, adapting to mobile viewports through Tailwind CSS breakpoint utilities and a collapsible hamburger navigation menu.

All asynchronous data fetching is managed through TanStack React Query~v5.
A \texttt{QueryClient} with a default stale time of five minutes is configured at the application root.
The Vite development server proxies all \texttt{/api} requests to the Django backend, configurable via the \texttt{VITE\_PROXY\_TARGET} environment variable.

Code quality tooling configured as part of this contribution includes ESLint, Prettier, and Husky pre-commit hooks that enforce linting and formatting before every commit.

The architecture deliberately avoids a global state management library.
All shared state is either query cache (React Query) or local component state (\texttt{useState}), which keeps the data flow unidirectional and makes individual components straightforward to test in isolation with Jest and Testing Library.

The analytics dashboard (\texttt{CatalogAnalyticsPage}) is populated from the \texttt{/api/sources/analytics/} endpoint and presents four visualisation panels:
\begin{enumerate}
    \item \textbf{Catalog comparison} — bar chart of source counts per catalog.
    \item \textbf{Detection trends} — time series of newly ingested sources by month.
    \item \textbf{Significance histogram} — log-spaced bins over the range $[10^0, 10^4]$ (15 bins), with the 95th-percentile significance value annotated.
    \item \textbf{Detectability breakdown} — stacked bar or pie chart showing the proportion of sources in each detectability tier (low, medium, high).
\end{enumerate}

\begin{figure}[h]
    \centering
    \includegraphics[width=0.8\linewidth]{skymap.png}
    \caption{Sky distribution of ingested sources colour-coded by primary catalog (Fermi-LAT, HAWC, LHAASO, TeVCat), shown in Galactic coordinates with an Aitoff projection, as rendered by the analytics sky map view built into the application frontend.}
    \label{fig:catalog_aggregation}
\end{figure}
He also focused on integrating MAGIC Source Simulator (MSS) detectability predictions into the analytics dashboard and visualizing source spectral energy distributions (SEDs) within the web application frontend.

The SED visualization component renders simulated spectral bins for sources where MAGIC detectability predictions are available, allowing astronomers to inspect the predicted gamma-ray flux as a function of observational energy bands.

\begin{figure}[h]
    \centering
    \includegraphics[width=1\linewidth]{image.png}
    \caption{Spectral energy distribution for 4FGL~J0532.6+0732 simulated by MAGIC Source Simulator for a 20-hour observation showing near-complete detection (12/13 bins) across 100~GeV--10~TeV with 100\% detection probability.}
    \label{fig:magic_sed_example}
\end{figure}

\subsection{Data normalization and schema handling}

The backend MSS integration exposes a \texttt{/api/sources/\{id\}/magic\_simulation/} endpoint that returns a JSON response containing an array of spectral points, each recording the simulated flux at a discrete energy bin alongside detection metadata.
To accommodate backward compatibility during schema evolution, the frontend normalisation layer accepts multiple naming conventions for the same physical quantity:
\begin{itemize}
    \item Integrated flux: \texttt{flux\_sed\_tev\_cm2\_s}
    \item Flux uncertainty: \texttt{flux\_sed\_error}
    \item Energy: \texttt{energy\_gev}
    \item Detection flag: \texttt{is\_detected}
    \item Statistical significance: \texttt{significance\_sigma}
\end{itemize}

Each incoming spectral point is normalized into an internal representation containing seven fields: energy, flux, fluxLow, fluxHigh, significance, isDetected, and an index. The asymmetric error bounds are computed as:
\begin{equation}
    \mathrm{fluxLow} = \mathrm{flux} - \Delta\mathrm{flux}, \quad \mathrm{fluxHigh} = \mathrm{flux} + \Delta\mathrm{flux}
\end{equation}
where $\Delta\mathrm{flux}$ is the reported flux uncertainty.
Points with non-finite or non-positive energy and flux values are filtered to ensure compatibility with logarithmic axis scaling.

\subsection{Chart rendering and log-space transformation}

The SED is visualized using Recharts, a declarative charting library built on React and D3. Both the energy (x-axis) and flux (y-axis) are rendered on logarithmic scales with the configuration \texttt{scale="log"} and \texttt{domain=["auto", "auto"]}, allowing Recharts to infer the axis bounds from the filtered point set.
Unlike some plotting libraries that require pre-logarithmic transformation of data, Recharts expects raw (linear-space) values and applies the log transformation internally; this ensures that error bar rendering and tick calculations remain numerically stable.

The chart displays two distinct series:
\begin{enumerate}
    \item \textbf{Detected points} — plotted as a scatter series with circular markers (colour \texttt{\#3266ad}), representing spectral bins where the simulated detection significance exceeds a threshold.
    \item \textbf{Upper limits} — plotted with downward-pointing triangle markers (colour \texttt{\#888780}), representing bins where no significant detection is predicted (e.g.\ low significance or low flux regimes).
\end{enumerate}

Axis tick positions are generated adaptively using a custom tick builder that identifies decade boundaries and intermediate points. The energy axis tick formatter produces labels in the form $10^n$ (using Unicode superscript characters) when the tick value coincides with a power of ten, and scientific notation otherwise. Intermediate ticks at $10^{n+0.5}$ (e.g.\ the geometric mean between decades) are included to reduce label density while maintaining readability across multiple orders of magnitude.

\subsection{Error bar visualization}

Asymmetric error bars are rendered using Recharts' \texttt{<ErrorBar>} component with a custom accessor function:
\begin{equation}
    \mathrm{errorBarAccessor}(d) = [\mathrm{fluxLow}, \mathrm{fluxHigh}]
\end{equation}
This accessor extracts the pre-computed bounds for each point. Because the y-axis uses logarithmic scaling, both $\mathrm{fluxLow}$ and $\mathrm{fluxHigh}$ must be strictly positive ($> 0$); points where the lower error bound is non-positive are not plotted. This represents a known edge case: if the MSS simulation predicts a low flux with a large relative uncertainty, the computed $\mathrm{fluxLow}$ may be negative or zero and cannot be displayed on a log scale. Recommended future mitigation includes clamping $\mathrm{fluxLow}$ to a small positive value (e.g.\ $\mathrm{flux} \times 10^{-30}$) or rendering a truncated arrow marker to indicate a lower bound below the axis range.

\subsection{Interactive tooltip and value extraction}

A custom \texttt{SedTooltip} component intercepts hover events and displays point-specific information. The Recharts tooltip payload contains complex nested structures with unknown types at runtime.

This avoids \texttt{any} type casts and provides graceful fallback to \texttt{NaN} for unparseable values. The tooltip displays:
\begin{itemize}
    \item Energy (formatted with human-readable units: GeV for $E < 1000$\,GeV, TeV for $E \ge 1000$\,GeV)
    \item Flux (scientific notation, two decimal places)
    \item Flux low and flux high bounds (scientific notation)
    \item Significance $\sigma$ (fixed-point, one decimal place)
\end{itemize}

The tooltip is styled with a semi-transparent dark card background and high-contrast foreground text to ensure readability over the chart plot area.

\begin{figure}[h]
    \centering
    \includegraphics[width=0.8\textwidth]{figures/placeholder_detectability.png}
    \caption{Distribution of the composite detectability score $S$ across all ingested sources, shown as a histogram with tier boundaries at $S=40$ (low/medium) and $S=70$ (medium/high) indicated by dashed vertical lines. The inset pie chart shows the fraction of sources in each tier, broken down by source class (PSR, AGN, SNR, unassociated).}
    \label{fig:detectability}
\end{figure}

\section{Main interface -- Mateusz Jaworski}

Mateusz Jaworski implemented the main application interface through which users explore and analyse the aggregated catalog data.
He was also responsible for the application routing and internationalisation.
Routing is handled by React Router~v7 in \texttt{BrowserRouter} mode; all routes share an \texttt{AppLayout} wrapper that renders the \texttt{NavigationHeader} and an \texttt{<Outlet>} for the active page, covering the paths \texttt{/}, \texttt{/universe-map}, \texttt{/source-analytics}, \texttt{/sources}, and \texttt{/sources/:id}.
Internationalisation is implemented with i18next~v26; both English and Polish translations are provided in \texttt{src/locales/en.json} and \texttt{src/locales/pl.json}, with the active language persisted in \texttt{localStorage} and toggled via the \texttt{LanguageSwitcher} component.

Mateusz Jaworski also implemented the custom React Query hooks that encapsulate all data fetching in the application: \texttt{useSourcesData} for the paginated and filtered sources list, \texttt{useSourceDetail} for individual source records, and \texttt{useSourceCatalogEntries} for the catalog entries of a given source.
Each hook surfaces \texttt{isLoading}, \texttt{isError}, and \texttt{data} states to consuming components, keeping data-fetching and presentation concerns cleanly separated from page rendering logic.

The sources list page (\texttt{SourcesPage}) fetches data from the \texttt{/api/sources/filter/} endpoint with server-side pagination and supports the following filter dimensions:

\begin{enumerate}
    \item \textbf{Catalog} — restrict results to one or more of \texttt{FERMI}, \texttt{HAWC}, \texttt{LHAASO}, \texttt{TEVCAT}.
    \item \textbf{Source class} — free-text match against the catalog metadata \texttt{source\_class} field (e.g.\ \texttt{PSR}, \texttt{AGN}, \texttt{SNR}).
    \item \textbf{Coordinate bounds} — right ascension $\alpha \in [0^{\circ}, 360^{\circ}]$ and declination $\delta \in [-90^{\circ}, +90^{\circ}]$ range sliders.
    \item \textbf{Detection significance} — minimum and maximum significance bounds.
    \item \textbf{Flux} — minimum and maximum integrated flux above 1\,GeV.
    \item \textbf{Minimum catalog count} — show only sources appearing in at least $n$ catalogs.
\end{enumerate}

Active filters are displayed as removable badges above the table, allowing users to inspect and clear individual constraints without navigating away.
Pagination state and filter values are managed with \texttt{useState} hooks local to the page component; changes to any filter reset the page index to one to avoid stale result windows.

The source detail page (\texttt{SourceDetailPage}) renders the unified \texttt{Source} record — name, equatorial coordinates, primary catalog, discovery date — followed by a table of all associated \texttt{CatalogEntry} records.
Each catalog entry row displays the original catalog name, the detection method, cross-match confidence, and the full metadata payload rendered as a key–value list, allowing domain experts to inspect native catalog quantities without loss of information.

All data-fetching in the main interface uses these custom React Query hooks, so loading and error states are handled consistently across all views.
The implementation keeps presentation components free of business logic: filter construction, parameter serialisation, and page-reset side effects are localised to the hook layer, making the page components straightforward to test and extend with additional filter dimensions in future releases.

\section{Skymap -- Krzysztof Wilk}

Krzysztof Wilk implemented both skymap views in the application: a 2D projected map rendered on an HTML \texttt{<canvas>} element using D3, and an interactive 3D celestial sphere built with React Three Fiber and Three.js.

The 2D sky map component (\texttt{CatalogSkyMap}) supports four cartographic projections: Aitoff, Mollweide, equirectangular, and orthographic.
The active projection is selected  at runtime from a dropdown control and applied via the corresponding \texttt{d3-geo} projection function.
Sources can be displayed in either the equatorial coordinate system $(\alpha, \delta)$ or the Galactic system $(l, b)$.
The conversion from equatorial to Galactic coordinates is implemented analytically using the standard transformation defined by the J2000 North Galactic Pole at $(\alpha_{\text{NGP}}, \delta_{\text{NGP}}) = (192.859\,48^{\circ},\; 27.128\,25^{\circ})$ and the ascending node offset $l_{\Omega} = 32.931\,92^{\circ}$:

\begin{equation}
    \sin b = \sin\delta\,\sin\delta_{\text{NGP}} + \cos\delta\,\cos\delta_{\text{NGP}}\,\cos(\alpha - \alpha_{\text{NGP}})
\end{equation}
\begin{equation}
    l = \arctan\!\frac{\cos\delta\,\sin(\alpha - \alpha_{\text{NGP}})}{\sin\delta\,\cos\delta_{\text{NGP}} - \cos\delta\,\sin\delta_{\text{NGP}}\,\cos(\alpha - \alpha_{\text{NGP}})} + l_{\Omega}
\end{equation}

Source points are drawn as filled circles whose radii are scaled with detection significance using a square-root scale:
\begin{equation}
    r = r_{\min} + (r_{\max} - r_{\min})\,\sqrt{\frac{\sigma}{\sigma_{\max}}}
    \qquad r_{\min} = 1.8\,\text{px},\quad r_{\max} = 6.4\,\text{px}
\end{equation}
Each source is coloured according to its primary catalog using a fixed colour palette.
The canvas is drawn on every frame update in a \texttt{useEffect} hook that iterates over all visible points and calls the Canvas 2D API directly, avoiding the overhead of DOM element creation per source.
Pan and zoom are implemented with the D3 zoom behaviour (scale extent $1\times$–$12\times$), with the projection transform updated on every zoom event.
Mouse hover interaction identifies the nearest projected point within a pixel-distance threshold and displays a tooltip with the source name, coordinates, and significance.
A cone-search filter accepts centre coordinates $(\alpha_0, \delta_0)$ and a radius $\rho$, restricting the rendered set to sources satisfying $\theta(\cdot, \cdot) \le \rho$.

The 3D view renders sources on a unit celestial sphere using Three.js geometry and React Three Fiber's declarative scene graph.
To maintain interactive frame rates over the full source set ($\sim$\,7000 objects), sources are spatially aggregated using a HEALPix-based hierarchical clustering scheme \citep{Gorski2005}.
The sky is subdivided into $N_{\text{side}} = 2^{\text{order}}$ HEALPix pixels in the NESTED indexing scheme; the \texttt{ang2pix} transformation maps a source's angular position $(\alpha, \delta)$ to its pixel index using two branches — an equatorial branch for $|z| = |\sin\delta| \le 2/3$ and a polar branch for $|z| > 2/3$ — mirroring the original Górski et al. algorithm \citep{Gorski2005}.
At coarser zoom levels, nearby sources are merged into cluster markers (\texttt{ClusterMarkers}) with a count badge; as the camera approaches a region the hierarchy is refined, revealing individual source points.

\begin{figure}[h!]
    \centering
    \includegraphics[width=0.8\textwidth]{figures/placeholder_skymap.png}
    \caption{2D sky map in Galactic Aitoff projection (left) showing all ingested sources colour-coded by catalog, with significance-scaled point radii. The 3D celestial sphere view (right) shows hierarchical clustering at an intermediate zoom level.}
    \label{fig:skymap}
\end{figure}

The two skymap views share the same underlying source data fetched via a dedicated API module that pages through the \texttt{/api/sources/} endpoint (up to $100 \times 200 = 20\,000$ items) and caches the result in the React Query store.
Future work could replace the paged fetch with a dedicated bulk-export endpoint returning compressed GeoJSON or HEALPix-indexed tiles to reduce initial load time for large catalog releases.

\clearpage
\begin{thebibliography}{99}
\bibitem[Abdollahi et al.(2020)]{2020ApJS..247...33A} Abdollahi, S., Acero, F., Ackermann, M., et al.\ 2020, ApJS, 247, 1, 33. doi:10.3847/1538-4365/ab6bcb
\bibitem[Li \& Ma(1983)]{LiMa1983} Li, T.-P., \& Ma, Y.-Q.\ 1983, ApJ, 272, 317. doi:10.1086/161295
\bibitem[Górski et al.(2005)]{Gorski2005} Górski, K.~M., Hivon, E., Banday, A.~J., et al.\ 2005, ApJ, 622, 759. doi:10.1086/427976
\end{thebibliography}

\end{document}
